\documentclass[12pt]{article}
\usepackage[utf8]{inputenc}
\usepackage[T1]{fontenc}
\usepackage{hyperref}
\title{Baduy Language Resources for K-12 Curriculum Development}
\author{Compiled by Claude Sonnet 5.5 (claude-sonnet-5-5) for Andr\'as Kornai}
\date{}
\begin{document}
\maketitle

\section{Introduction}
Baduy (ISO 639-3 bac) is spoken by the Baduy (Kanekes) people of Banten, Indonesia. It is conventionally treated as a Sundanese dialect, though its distinct vocabulary leads many to count it as a separate language. It has no official writing system. Research is relatively limited and no dictionary or grammar turned up in the searches.

\section{Primary Educational Resources}
None found.

\section{Linguistic and Reference Materials}
\begin{itemize}
\item Wikipedia entry on the Baduy language. \url{https://en.wikipedia.org/wiki/Baduy_language}
\end{itemize}

\section{Religious and Community Resources}
\begin{itemize}
\item Wikipedia entry on the Baduy people. \url{https://en.wikipedia.org/wiki/Baduy_people}
\end{itemize}

\section{Digital and Audio Resources}
\begin{itemize}
\item PARADISEC item BAC01-024, a recorded narrative from the Baduy Dalam collection. \url{https://catalog.paradisec.org.au/repository/BAC01/024}
\item Research Data Australia record, ``Documenting endangered Indonesian Baduy Dalam''. \url{https://researchdata.edu.au/documenting-endangered-indonesian-baduy-dalam/2369835}
\end{itemize}

\section{Curriculum Development Context}
The PARADISEC recordings are the main primary data found. A curriculum would depend on new documentation and on decisions about orthography with the community.

\end{document}